\begin{summarybox}
\textbf{【一句话记忆】椭圆平方拆成“切点乘变量”，前提先看点在何处。}

\textbf{1. 最核心公式：}
\[
 \boxed{\frac{x_0x}{a^2}+\frac{y_0y}{b^2}=1.}
\]
若 $P(x_0,y_0)$ 在椭圆上，它是 $P$ 处切线；若给的是椭圆外点 $M(x_1,y_1)$，换成 $x_1,y_1$ 得到的是\textbf{两切点的连线（极线）}，不是经过 $M$ 的切线。

\textbf{2. 最快证明（平方和一行）：}

设 $Q$ 同时在椭圆及上述直线上，则
\[
 \frac{(x-x_0)^2}{a^2}+\frac{(y-y_0)^2}{b^2}
 =1+1-2\times1=0
 \quad\Longrightarrow\quad Q=P.
\]
故此直线与椭圆仅在 $P$ 处相切。\textbf{无需导数、判别式、斜率分类。}

\textbf{3. 最直观理解：}令 $u=x/a,v=y/b$，椭圆变成圆 $u^2+v^2=1$；\textbf{圆的切线垂直半径}，变回原坐标就是切点式。

\textbf{4. 两个高考常用配套公式：}
\[
 \boldsymbol n=\left(\frac{x_0}{a^2},\frac{y_0}{b^2}\right),
 \qquad k=-\frac{b^2x_0}{a^2y_0}\ (y_0\ne0),
\]
\[
 Ax+By+C=0\text{ 与椭圆相切}\iff
 \boxed{C^2=a^2A^2+b^2B^2}\quad((A,B)\ne(0,0)).
\]

\textbf{5. 实战流程：}先检查点是否在椭圆上（$S=1$）、外（$S>1$）或内（$S<1$）；再按需求写\textbf{切线、切点弦或相切判定}。遇到坐标轴端点优先用切点式，避免斜率除零。

\textbf{6. 最易错：}外点的极线不穿过外点；椭圆法线通常不等于原点到切点的连线；“切点式”不要与“点幂定理”混淆。
\end{summarybox}
